\Needspace{8\baselineskip}
\Article{212}
\textsc{Problem.} \textit{Find all representations of a given number $M$ by a given form $axx+2bxy+cyy$ with determinant $hh$.}

We can solve this problem from the principles of Article~168 in exactly the same way as we showed above (Articles~180, 181, 205) for forms with negative determinant and positive nonsquare determinant. Since this involves no difficulty, it would be unnecessary to repeat it here. It is useful, however, to derive the solution from another principle specific to the present case.

Putting, as in Articles~206, 208,
\[
\begin{gathered}
h-b:a=c:-(h+b)=\beta:\delta\\
\frac{h-b}{\beta}=\frac a\delta=f;\quad
\frac c\beta=\frac{-h-b}{\delta}=g
\end{gathered}
\]
we can easily prove that the given form is the product of the factors $\delta x-\beta y$ and $fx-gy$. So every representation of $M$ by the given form gives a factorization of $M$ into two factors. If all divisors of $M$ are $d,d',d''$ etc. (including $1$ and $M$, and taking each \textit{twice}, both positively and negatively), we clearly obtain all representations of $M$ by putting in turn
\[
\begin{aligned}
\delta x-\beta y&=d,&\quad fx-gy&=\frac M d\\
\delta x-\beta y&=d',& fx-gy&=\frac M{d'}\quad\text{etc.}
\end{aligned}
\]
finding the values of $x$, $y$ from these equations, and discarding the representations in which $x$
\SourceLeaf{211}
or $y$ has a fractional value. From the first two equations we clearly get
\[
x=\tfrac{\beta M-gdd}{(\beta f-\delta g)d}\quad\text{and}\quad
y=\tfrac{\delta M-fdd}{(\beta f-\delta g)d}
\]
These values are always \textit{determinate}, since $\beta f-\delta g=2h$, so the numerator is certainly not $=0$. --- We could have solved the other problems from the same principle, namely that every form with square determinant factors into two factors. But here too we preferred to use a method analogous to the one given above for forms with nonsquare determinant.

\textit{Example.} We seek all representations of $12$ by $3xx+4xy-7yy$. This factors into $x-y$ and $3x+7y$. All divisors of $12$ are $\pm1,2,3,4,6,12$. Putting $x-y=1$, $3x+7y=12$ gives $x=\frac{19}{10}$, $y=\frac9{10}$, which we must discard because they are fractional. In the same way the divisors $-1,\pm3,\pm4,\pm6,\pm12$ give unusable values. But the divisor $+2$ gives $x=2,y=0$, and $-2$ gives $x=-2$, $y=0$. So there are no representations beyond these two.

This method cannot be used if $M=0$. In that case all values of $x,y$ must clearly satisfy either $\delta x-\beta y=0$ or $fx-gy=0$. All solutions of the first equation are contained in the formula $x=\beta z$, $y=\delta z$, where $z$ ranges over all integers (since, as assumed, $\beta$, $\delta$ are relatively prime). Similarly, putting the greatest common divisor of $f,g$ $=m$, all solutions of the second equation are given by $x=\frac{gz}{m},y=\frac{hz}{m}$. These two general formulas therefore contain all representations of $M$ in this case.

\begin{center}
\(\begin{matrix} *&\qquad&*\\[-0.8ex]&*&\end{matrix}\)
\end{center}

In the preceding discussion we have set out everything needed to recognize equivalence, find all transformations of forms, and find all representations of given numbers by given forms, so that there seems to be nothing more to ask for. All that remains is to show, given two forms that cannot be equivalent because their \textit{determinants are unequal}, how to decide whether one is contained in the other and, if so, how to find all transformations of the former into the latter.

